\begin{abstract}
This paper studies whether pre-announcement trading pressure helps explain cross-sectional variation in post-earnings-announcement drift (PEAD). While the PEAD literature documents that stock returns continue to drift following earnings surprises \citep{BernardThomas1989}, it typically measures news using accounting-based signals and abstracts from the microstructure environment in which prices adjust. In parallel, market microstructure models show that order flow conveys private information and that trades can have either persistent (informational) or transient (liquidity-driven) price impacts \citep{Kyle1985,Hasbrouck1991}.

We bridge these literatures by constructing a continuous, event-level \emph{pressure score} from intraday TAQ data in the pre-announcement window, designed to proxy for the relative intensity of information-driven versus liquidity-driven trading. We test whether this microstructure-based score predicts heterogeneity in post-earnings returns conditional on standardized unexpected earnings (SUE).

To complement this economically motivated construction, we apply principal component analysis (PCA) to the microstructure feature space. Rather than imposing a priori structure, we use PCA to extract latent factors and map them to economically interpretable dimensions of trading behavior. We show that these PCA-derived factors—particularly those capturing persistent and directional order flow—have direct predictive content for post-earnings returns.

Our findings indicate that stronger information-related trading pressure prior to earnings announcements is associated with attenuated subsequent drift, consistent with partial information incorporation before the event. Overall, the results highlight that return dynamics around earnings announcements depend not only on the magnitude of information, but also on the timing and mechanism through which that information is incorporated into prices.
\end{abstract}
